\chapter{Calcul du gradient}
\label{ch:gradient}

\begin{lecture}
Le coût est un nombre qui dépend des coefficients. Pour le faire baisser, il faut
savoir comment il varie quand chaque coefficient varie : c'est sa dérivée
partielle. Le calcul se décompose en deux dérivées enchaînées, et le résultat se
réduit à l'erreur commise, multipliée par la note.
\end{lecture}

\section{Composition à dériver}

Le coefficient $w_j$ n'apparaît pas directement dans $J$. Il agit par étapes :
\[
  w_j \;\longrightarrow\; z_i=\sum_j w_jx_{ij}
  \;\longrightarrow\; p_i=\sigma(z_i)
  \;\longrightarrow\; \ell_i \;\longrightarrow\; J .
\]
La règle de dérivation des fonctions composées traverse cette suite : chaque
facteur se calcule séparément, puis les facteurs se multiplient.

\begin{figure}[htbp]
\centering
\begin{tikzpicture}[
  bx/.style={draw=ink!30,rounded corners=2.5pt,align=center,font=\small,
    inner sep=5pt,minimum width=17mm,minimum height=10mm},
  arr/.style={-{Latex[length=2.2mm]},thick,ink!60},
  der/.style={font=\scriptsize,text=accent,align=center}]
\node[bx,fill=gridgray!20] (w) {$w_j$};
\node[bx,fill=accentlight,right=14mm of w] (z) {$z_i$};
\node[bx,fill=accentlight,right=14mm of z] (p) {$p_i$};
\node[bx,fill=warninglight,right=14mm of p] (l) {$\ell_i$};
\node[bx,fill=rulelight,right=14mm of l] (j) {$J$};
\draw[arr] (w) -- node[der,above=1mm]{$x_{ij}$} (z);
\draw[arr] (z) -- node[der,above=1mm]{$\sigma$} (p);
\draw[arr] (p) -- (l);
\draw[arr] (l) -- node[der,above=1mm]{moyenne} (j);
\draw[{Latex[length=2.2mm]}-,thick,warning!70]
  ($(w.south)+(0,-4mm)$) -- ($(z.south)+(0,-4mm)$);
\node[der,text=warning,below=6mm of w,xshift=9.5mm,anchor=north]
  {$\dfrac{\partial z_i}{\partial w_j}=x_{ij}$};
\draw[{Latex[length=2.2mm]}-,thick,warning!70]
  ($(z.south)+(0,-4mm)$) -- ($(l.south)+(0,-4mm)$);
\node[der,text=warning,below=6mm of p,anchor=north]
  {$\dfrac{\partial \ell_i}{\partial z_i}=p_i-y_i$};
\end{tikzpicture}
\caption{Le coefficient n'agit sur le coût que par cette chaîne. Les deux
dérivées encadrées se multiplient, et la moyenne sur les élèves donne
$\partial J/\partial w_j$. Le passage par $p_i$ disparaît du résultat : la
propriété $\sigma'=\sigma(1-\sigma)$ fait que les deux dérivées intermédiaires se
simplifient en un seul facteur $p_i-y_i$.}
\label{fig:chaine}
\end{figure}
\FloatBarrier

\section{Première dérivée : le coût d'un élève par rapport au score}

Posons, pour un élève dont on omet l'indice,
\[
  \ell=-\bigl[y\log\sigma(z)+(1-y)\log(1-\sigma(z))\bigr].
\]

\begin{align}
  \frac{\partial \ell}{\partial z}
  &= -\left[y\cdot\frac{1}{\sigma(z)}\cdot\sigma'(z)
     +(1-y)\cdot\frac{1}{1-\sigma(z)}\cdot\bigl(-\sigma'(z)\bigr)\right]
    &&\text{chaîne sur } \log \notag\\[4pt]
  &= -\left[\frac{y\,\sigma'(z)}{\sigma(z)}
     -\frac{(1-y)\,\sigma'(z)}{1-\sigma(z)}\right]
    &&\text{on range} \notag\\[4pt]
  &= -\left[\frac{y\,\sigma(z)(1-\sigma(z))}{\sigma(z)}
     -\frac{(1-y)\,\sigma(z)(1-\sigma(z))}{1-\sigma(z)}\right]
    &&\sigma'=\sigma(1-\sigma) \notag\\[4pt]
  &= -\bigl[y\bigl(1-\sigma(z)\bigr)-(1-y)\sigma(z)\bigr]
    &&\text{on simplifie chaque fraction} \notag\\[4pt]
  &= -\bigl[y-y\,\sigma(z)-\sigma(z)+y\,\sigma(z)\bigr]
    &&\text{on développe} \notag\\[4pt]
  &= -\bigl[y-\sigma(z)\bigr]
    &&\text{les } y\,\sigma(z) \text{ s'annulent} \notag\\[4pt]
  &= \sigma(z)-y \label{eq:dl-dz}
\end{align}

Le troisième passage repose entièrement sur la propriété
$\sigma'=\sigma(1-\sigma)$ (ch.~\ref{ch:sigmoide}) : c'est elle qui fait
disparaître les dénominateurs.

\begin{resultat}[Résidu]
\[
  \frac{\partial \ell}{\partial z}=\sigma(z)-y=p-y .
\]
La pente du coût par rapport au score est simplement l'écart entre ce que le
modèle annonce et ce qui est observé. Un élève \Gryf ($y=1$) auquel le
modèle accorde $0{,}1$ donne $-0{,}9$ ; le même, déjà classé à $0{,}99$, donne
$-0{,}01$. La correction issue d'un élève correctement classé est donc
négligeable.
\end{resultat}

\section{Seconde dérivée : le score par rapport à un coefficient}

\begin{align}
  \frac{\partial z_i}{\partial w_j}
  &= \frac{\partial}{\partial w_j}\sum_{k=0}^{n}w_kx_{ik} \notag\\[2pt]
  &= \sum_{k=0}^{n}\frac{\partial (w_kx_{ik})}{\partial w_j}
    &&\text{la dérivée entre dans la somme} \notag\\[2pt]
  &= x_{ij}
    &&\text{dérivée partielle}
\end{align}

Dans la somme, un seul terme contient $w_j$, c'est $w_jx_{ij}$, et sa dérivée par
rapport à $w_j$ vaut $x_{ij}$. Les autres termes $w_kx_{ik}$, avec $k\neq j$, sont
des constantes du point de vue de cette dérivée, qui est donc nulle.

\section{Assemblage}

\begin{align}
  \frac{\partial J}{\partial w_j}
  &= \frac{\partial}{\partial w_j}\left[\frac{1}{m}\sum_{i=1}^{m}\ell_i\right] \notag\\[2pt]
  &= \frac{1}{m}\sum_{i=1}^{m}\frac{\partial \ell_i}{\partial w_j}
    &&\text{la dérivée entre, } 1/m \text{ sort} \notag\\[2pt]
  &= \frac{1}{m}\sum_{i=1}^{m}
     \frac{\partial \ell_i}{\partial z_i}\cdot\frac{\partial z_i}{\partial w_j}
    &&\text{règle de chaîne} \notag\\[2pt]
  &= \frac{1}{m}\sum_{i=1}^{m}\bigl(\sigma(z_i)-y_i\bigr)\,x_{ij}
    &&\text{les deux dérivées calculées ci-dessus}
\end{align}

\begin{resultat}[Gradient]
\[
  \boxed{\;\frac{\partial J}{\partial w_j}
  =\ann{rulecolor}{\frac{1}{m}\sum_{i=1}^{m}}{\annl{moyenne}}
   \;\ann{warning}{\bigl(p_i-y_i\bigr)}{\annl{erreur}}
   \;\ann{accent}{x_{ij}}{\annl{note}}\;}
\]
C'est la formule du sujet, $\frac{1}{m}\sum(h_\theta(x^i)-y^i)x^i_j$, obtenue sans
hypothèse supplémentaire. Le gradient $\nabla J$ est le vecteur de ces $n+1$
dérivées partielles, une par coefficient.
\end{resultat}

\section{Facteurs du gradient}

Chaque élève contribue au gradient du coefficient $j$ par le produit de deux
facteurs : son erreur $p_i-y_i$, et sa note $x_{ij}$.

Le premier facteur donne le sens et l'intensité de la correction. Il est négatif
quand le modèle sous-estime un \Gryf, positif quand il surestime un autre
élève, et proche de zéro dès que la prédiction est correcte.

Le second répartit la correction entre les coefficients. Un élève dont la note
d'\Herb est extrême pèse fortement sur $w_1$, tandis qu'un élève de note
moyenne, dont la valeur standardisée est proche de zéro, n'y contribue presque
pas.

\begin{vigilance}[Cas du biais]
Pour $j=0$, la note $x_{i0}$ vaut $1$ pour tout le monde. La formule devient
\[
  \frac{\partial J}{\partial w_0}=\frac{1}{m}\sum_{i=1}^{m}(p_i-y_i),
\]
soit la moyenne des erreurs. Le biais s'ajuste donc sur la fréquence globale : il
se déplace jusqu'à ce que la proportion moyenne annoncée égale la proportion
réelle de \Gryf, ce que la première itération vérifie
(ch.~\ref{ch:descente}).
\end{vigilance}
